% answers.tex —— 2005 年·普通高等学校招生考试·重庆卷·文科 · 参考答案与详细解析片段
% 由 tools/generate_projects.py 依来源总账自 2005/2005_数学真题_重庆文科_试题.md + 2005/2005_数学真题_重庆文科_解析正文.md + 2005/2005_数学真题_重庆文科_答案要点.md 生成，勿手改（改源 MD 后重跑）

\begin{answer}
\gksolhead{1}{}

\textbf{题目}：

圆 $(x + 2)^2 + y^2 = 5$ 关于原点 $(0, 0)$ 对称的圆的方程为\mbox{（　　）}

\begin{gkopts}(2)
\task A．$(x - 2)^2 + y^2 = 5$
\task B．$x^2 + (y - 2)^2 = 5$
\task C．$(x + 2)^2 + (y + 2)^2 = 5$
\task D．$x^2 + (y + 2)^2 = 5$
\end{gkopts}

\Needspace{4\baselineskip}
\textbf{解析}

原圆方程为 $(x + 2)^2 + y^2 = 5$，其圆心坐标为 $(-2, 0)$，半径为 $\sqrt{5}$．
圆心 $(-2, 0)$ 关于原点 $(0, 0)$ 的对称点坐标为 $(2, 0)$．
因为对称后的圆半径保持不变，仍为 $\sqrt{5}$，
所以所求圆的方程为 $(x - 2)^2 + y^2 = 5$．
故选 A．

\gkanswer{答案：A}

\end{answer}
\begin{answer}
\gksolhead{2}{}

\textbf{题目}：

$\left(\cos\dfrac{\pi}{12} - \sin\dfrac{\pi}{12}\right)\left(\cos\dfrac{\pi}{12} + \sin\dfrac{\pi}{12}\right) =$\mbox{（　　）}

\begin{gkopts}(2)
\task A．$-\dfrac{\sqrt{3}}{2}$
\task B．$-\dfrac{1}{2}$
\task C．$\dfrac{1}{2}$
\task D．$\dfrac{\sqrt{3}}{2}$
\end{gkopts}

\Needspace{4\baselineskip}
\textbf{解析}

\Needspace{5\baselineskip}
利用平方差公式与二倍角余弦公式展开：
\nopagebreak[4]
\[
\left(\cos\frac{\pi}{12} - \sin\frac{\pi}{12}\right)\left(\cos\frac{\pi}{12} + \sin\frac{\pi}{12}\right) = \cos^2\frac{\pi}{12} - \sin^2\frac{\pi}{12}.
\]
\Needspace{5\baselineskip}
由二倍角公式 $\cos 2\theta = \cos^2\theta - \sin^2\theta$，得：
\nopagebreak[4]
\[
\cos^2\frac{\pi}{12} - \sin^2\frac{\pi}{12} = \cos\left(2 \times \frac{\pi}{12}\right) = \cos\frac{\pi}{6} = \frac{\sqrt{3}}{2}.
\]
故选 D．

\gkanswer{答案：D}

\end{answer}
\begin{answer}
\gksolhead{3}{}

\textbf{题目}：

若函数 $f(x)$ 是定义在 $\mathbf{R}$ 上的偶函数，在 $(-\infty, 0]$ 上是减函数，且 $f(2) = 0$，则使得 $f(x) < 0$ 的 $x$ 的取值范围是\mbox{（　　）}

\begin{gkopts}(2)
\task A．$(-\infty, 2)$
\task B．$(2, +\infty)$
\task C．$(-\infty, -2) \cup (2, +\infty)$
\task D．$(-2, 2)$
\end{gkopts}

\Needspace{4\baselineskip}
\textbf{解析}

因为 $f(x)$ 是定义在 $\mathbf{R}$ 上的偶函数，其图象关于 $y$ 轴对称．
又 $f(x)$ 在 $(-\infty, 0]$ 上是减函数，由偶函数的对称性可知，$f(x)$ 在 $[0, +\infty)$ 上是增函数．
已知 $f(2) = 0$，由偶函数性质可得 $f(-2) = f(2) = 0$．
在区间 $[0, +\infty)$ 上：
当 $0 \leqslant x < 2$ 时，$f(x) < f(2) = 0$；
由对称性，当 $-2 < x \leqslant 0$ 时，$f(x) = f(-x) < f(2) = 0$．
综上，使得 $f(x) < 0$ 的 $x$ 的取值范围是 $(-2, 2)$．
故选 D．

\gkanswer{答案：D}

\end{answer}
\begin{answer}
\gksolhead{4}{}

\textbf{题目}：

设向量 $\vec{a} = (-1, 2), \vec{b} = (2, -1)$，则 $(\vec{a} \cdot \vec{b})(\vec{a} + \vec{b})$ 等于\mbox{（　　）}

\begin{gkopts}(4)
\task A．$(1, 1)$
\task B．$(-4, -4)$
\task C．$-4$
\task D．$(-2, -2)$
\end{gkopts}

\Needspace{4\baselineskip}
\textbf{解析}

已知 $\vec{a} = (-1, 2)$，$\vec{b} = (2, -1)$．
\Needspace{5\baselineskip}
先求数量积：
\nopagebreak[4]
\[
\vec{a} \cdot \vec{b} = (-1) \times 2 + 2 \times (-1) = -2 - 2 = -4.
\]
\Needspace{5\baselineskip}
再求向量和：
\nopagebreak[4]
\[
\vec{a} + \vec{b} = (-1 + 2, 2 - 1) = (1, 1).
\]
\Needspace{5\baselineskip}
因此：
\nopagebreak[4]
\[
(\vec{a} \cdot \vec{b})(\vec{a} + \vec{b}) = -4(1, 1) = (-4, -4).
\]
故选 B．

\gkanswer{答案：B}

\end{answer}
\begin{answer}
\gksolhead{5}{}

\textbf{题目}：

不等式组 $\begin{cases} |x - 2| < 2, \\ \log_2(x^2 - 1) > 1 \end{cases}$ 的解集为\mbox{（　　）}

\begin{gkopts}(4)
\task A．$(0, \sqrt{3})$
\task B．$(\sqrt{3}, 2)$
\task C．$(\sqrt{3}, 4)$
\task D．$(2, 4)$
\end{gkopts}

\Needspace{4\baselineskip}
\textbf{解析}

\Needspace{5\baselineskip}
解不等式 $|x - 2| < 2$：
\nopagebreak[4]
\[
-2 < x - 2 < 2 \iff 0 < x < 4.
\]
解对数不等式 $\log_2(x^2 - 1) > 1$：
\Needspace{5\baselineskip}
真数需满足：
\nopagebreak[4]
\[
x^2 - 1 > 2^1 = 2 \iff x^2 > 3 \iff x > \sqrt{3} \text{ \text{或} } x < -\sqrt{3}.
\]
\Needspace{5\baselineskip}
取两个不等式解集的交集：
\nopagebreak[4]
\[
(0, 4) \cap \left((-\infty, -\sqrt{3}) \cup (\sqrt{3}, +\infty)\right) = (\sqrt{3}, 4).
\]
故选 C．

\gkanswer{答案：C}

\end{answer}
\begin{answer}
\gksolhead{6}{}

\textbf{题目}：

已知 $\alpha$、$\beta$ 均为锐角，若 $p : \sin\alpha < \sin(\alpha + \beta)$，$q : \alpha + \beta < \dfrac{\pi}{2}$，则 $p$ 是 $q$ 的\mbox{（　　）}

\begin{gkopts}(2)
\task A．充分而不必要条件
\task B．必要而不充分条件
\task C．充要条件
\task D．既不充分也不必要条件
\end{gkopts}

\Needspace{4\baselineskip}
\textbf{解析}

因为 $\alpha, \beta$ 均为锐角，所以 $\alpha \in \left(0, \dfrac{\pi}{2}\right)$，$\beta \in \left(0, \dfrac{\pi}{2}\right)$，从而 $\alpha < \alpha + \beta < \pi$．
若条件 $q$ 成立，即 $\alpha + \beta < \dfrac{\pi}{2}$，则 $0 < \alpha < \alpha + \beta < \dfrac{\pi}{2}$．
由于正弦函数 $y = \sin x$ 在 $\left(0, \dfrac{\pi}{2}\right)$ 上单调递增，故 $\sin\alpha < \sin(\alpha + \beta)$，即条件 $p$ 成立．
所以 $q \implies p$．
反之，若条件 $p$ 成立，取 $\alpha = \dfrac{\pi}{6}$，$\beta = \dfrac{\pi}{3}$，
此时 $\alpha + \beta = \dfrac{\pi}{2}$，
$\sin\alpha = \dfrac{1}{2}$，$\sin(\alpha + \beta) = \sin\dfrac{\pi}{2} = 1 > \dfrac{1}{2}$，满足条件 $p$；
但此时 $\alpha + \beta = \dfrac{\pi}{2}$ 不满足 $\alpha + \beta < \dfrac{\pi}{2}$（条件 $q$ 不成立）．
故由 $p$ 推不出 $q$．
因此 $p$ 是 $q$ 的必要而不充分条件．
故选 B．

\gkanswer{答案：B}

\end{answer}
\begin{answer}
\gksolhead{7}{}

\textbf{题目}：

对于不重合的两个平面 $\alpha$ 与 $\beta$，给定下列条件：

① 存在平面 $\gamma$，使得 $\alpha$、$\beta$ 都垂直于 $\gamma$；

② 存在平面 $\gamma$，使得 $\alpha$、$\beta$ 都平行于 $\gamma$；

③ 存在直线 $l \subset \alpha$，直线 $m \subset \alpha$，使得 $l \parallel m$；

④ 存在异面直线 $l$、$m$，使得 $l \parallel \alpha$，$l \parallel \beta$，$m \parallel \alpha$，$m \parallel \beta$．

其中，可以判定 $\alpha$ 与 $\beta$ 平行的条件有\mbox{（　　）}

\begin{gkopts}(4)
\task A．1 个
\task B．2 个
\task C．3 个
\task D．4 个
\end{gkopts}

\Needspace{4\baselineskip}
\textbf{解析}

逐一分析各个条件：
① 存在平面 $\gamma$，使得 $\alpha \perp \gamma$ 且 $\beta \perp \gamma$．
垂直于同一平面的两个平面可以平行，也可以相交（例如建筑相邻的两个墙面均垂直于地面，但二者相交），故不能判定 $\alpha \parallel \beta$；
② 存在平面 $\gamma$，使得 $\alpha \parallel \gamma$ 且 $\beta \parallel \gamma$．
根据空间平面平行的传递性，平行于同一平面的两个不同平面平行，故能判定 $\alpha \parallel \beta$；
③ 存在直线 $l \subset \alpha, m \subset \alpha$ 使得 $l \parallel m$．
此条件仅反映平面 $\alpha$ 内部有两条平行线，与平面 $\beta$ 无位置关联，故不能判定 $\alpha \parallel \beta$；
④ 存在异面直线 $l, m$，使得 $l \parallel \alpha, l \parallel \beta$ 且 $m \parallel \alpha, m \parallel \beta$．
在平面 $\alpha$ 内可以作直线 $l' \parallel l, m' \parallel m$．
由于 $l$ 与 $m$ 是异面直线，故 $l'$ 与 $m'$ 是相交直线．
又因 $l \parallel \beta, m \parallel \beta$，所以 $l' \parallel \beta, m' \parallel \beta$．
由面面平行的判定定理（一个平面内两条相交直线分别平行于另一个平面，则这两个平面平行），可得 $\alpha \parallel \beta$．能判定平行．
综上，可以判定 $\alpha$ 与 $\beta$ 平行的条件为 ② 和 ④，共有 2 个．
故选 B．

\gkanswer{答案：B}

\end{answer}
\begin{answer}
\gksolhead{8}{}

\textbf{题目}：

若 $(1 + 2x)^n$ 展开式中含 $x^3$ 项的系数等于含 $x$ 项的系数的 8 倍，则 $n$ 等于\mbox{（　　）}

\begin{gkopts}(4)
\task A．$5$
\task B．$7$
\task C．$9$
\task D．$11$
\end{gkopts}

\Needspace{4\baselineskip}
\textbf{解析}

\Needspace{5\baselineskip}
$(1 + 2x)^n$ 的二项展开式通项为：
\nopagebreak[4]
\[
T_{r+1} = \mathrm{C}_n^r (2x)^r = \mathrm{C}_n^r 2^r x^r.
\]
含 $x$ 项对应 $r = 1$，其系数为 $\mathrm{C}_n^1 \times 2 = 2n$．
\Needspace{5\baselineskip}
含 $x^3$ 项对应 $r = 3$，其系数为：
\nopagebreak[4]
\[
\mathrm{C}_n^3 \times 2^3 = 8 \times \frac{n(n-1)(n-2)}{6} = \frac{4}{3}n(n-1)(n-2).
\]
\Needspace{5\baselineskip}
依题意，含 $x^3$ 项的系数等于含 $x$ 项系数的 8 倍：
\nopagebreak[4]
\[
\frac{4}{3}n(n-1)(n-2) = 8 \times (2n) = 16n.
\]
\Needspace{5\baselineskip}
因为 $n \geqslant 3$，两边同时除以 $4n$，得：
\nopagebreak[4]
\[
\frac{1}{3}(n-1)(n-2) = 4 \implies (n-1)(n-2) = 12.
\]
因为 $n \in \mathbf{N}^*$，解得 $n - 1 = 4 \implies n = 5$．
故选 A．

\gkanswer{答案：A}

\end{answer}
\begin{answer}
\gksolhead{9}{}

\textbf{题目}：

若动点 $(x, y)$ 在曲线 $\dfrac{x^2}{4} + \dfrac{y^2}{b^2} = 1$（$b > 0$）上变化，则 $x^2 + 2y$ 的最大值为\mbox{（　　）}

\begin{gkopts}(1)
\task A．$\begin{cases} \dfrac{b^2}{4} + 4, & 0 < b < 4, \\ 2b, & b \geqslant 4 \end{cases}$
\task B．$\begin{cases} \dfrac{b^2}{4} + 4, & 0 < b < 2, \\ 2b, & b \geqslant 2 \end{cases}$
\task C．$\dfrac{b^2}{4} + 4$
\task D．$2b$
\end{gkopts}

\Needspace{4\baselineskip}
\textbf{解析}

\Needspace{5\baselineskip}
由曲线方程 $\dfrac{x^2}{4} + \dfrac{y^2}{b^2} = 1$（$b > 0$），解得：
\nopagebreak[4]
\[
x^2 = 4\left(1 - \frac{y^2}{b^2}\right), \quad y \in [-b, b].
\]
\Needspace{5\baselineskip}
代入 $x^2 + 2y$ 得：
\nopagebreak[4]
\[
f(y) = 4 - \frac{4}{b^2}y^2 + 2y = -\frac{4}{b^2}\left(y^2 - \frac{b^2}{2}y\right) + 4 = -\frac{4}{b^2}\left(y - \frac{b^2}{4}\right)^2 + \frac{b^2}{4} + 4.
\]
这是一个关于 $y$ 的二次函数，抛物线开口向下，对称轴为 $y = \dfrac{b^2}{4}$．
结合定义域 $y \in [-b, b]$：
\Needspace{5\baselineskip}
比较对称轴与区间上界 $b$ 的大小：
\nopagebreak[4]
\[
\frac{b^2}{4} \leqslant b \iff b(b - 4) \leqslant 0 \iff 0 < b \leqslant 4.
\]
\begin{itemize}
\item 当 $0 < b < 4$ 时，对称轴 $y = \dfrac{b^2}{4} \in [-b, b]$，二次函数在顶点处取得最大值：
\end{itemize}
\nopagebreak[4]
\[
f_{\max} = \frac{b^2}{4} + 4.
\]
\begin{itemize}
\item 当 $b \geqslant 4$ 时，对称轴 $y = \dfrac{b^2}{4} \geqslant b$，$f(y)$ 在 $[-b, b]$ 上单调递增，在右端点 $y = b$ 处取得最大值：
\end{itemize}
\nopagebreak[4]
\[
f_{\max} = f(b) = -\frac{4}{b^2}b^2 + 2b + 4 = 2b.
\]
\Needspace{5\baselineskip}
综上所述：
\nopagebreak[4]
\[
f_{\max} = \begin{cases} \dfrac{b^2}{4} + 4, & 0 < b < 4, \\ 2b, & b \geqslant 4. \end{cases}
\]
故选 A．

\gkanswer{答案：A}

\end{answer}
\begin{answer}
\gksolhead{10}{}

\textbf{题目}：

有一塔形几何体由若干个正方体构成，构成方式如图所示，上层正方体下底面的四个顶点是下层正方体上底面各连接中点，已知最底层正方体的棱长为 2，且该塔形的表面积（含最底层正方体的底面面积）超过 39，则该塔形中正方体的个数至少是\mbox{（　　）}
\nopagebreak[4]
\begin{center}\begin{tikzpicture}[x=1cm, y=1cm, line join=round, line cap=round, thick]
  \def\dx{0.45}
  \def\dy{0.38}

  % ================= Layer 1 (Bottom cube) =================
  \coordinate (A1) at (0, 0);
  \coordinate (B1) at (1.6, 0);
  \coordinate (C1) at ($(B1) + (\dx, \dy)$);
  \coordinate (D1) at ($(A1) + (\dx, \dy)$);

  \def\hOne{1.6}
  \coordinate (A1t) at ($(A1) + (0, \hOne)$);
  \coordinate (B1t) at ($(B1) + (0, \hOne)$);
  \coordinate (C1t) at ($(C1) + (0, \hOne)$);
  \coordinate (D1t) at ($(D1) + (0, \hOne)$);

  \filldraw[fill=white] (A1t) -- (B1t) -- (C1t) -- (D1t) -- cycle;
  \filldraw[fill=white] (A1) -- (B1) -- (B1t) -- (A1t) -- cycle;
  \filldraw[fill=white] (B1) -- (C1) -- (C1t) -- (B1t) -- cycle;

  % ================= Layer 2 (Middle cube) =================
  \coordinate (M1) at ($(A1t)!0.5!(B1t)$);
  \coordinate (M2) at ($(B1t)!0.5!(C1t)$);
  \coordinate (M3) at ($(C1t)!0.5!(D1t)$);
  \coordinate (M4) at ($(D1t)!0.5!(A1t)$);

  \def\hTwo{1.13}
  \coordinate (M1t) at ($(M1) + (0, \hTwo)$);
  \coordinate (M2t) at ($(M2) + (0, \hTwo)$);
  \coordinate (M3t) at ($(M3) + (0, \hTwo)$);
  \coordinate (M4t) at ($(M4) + (0, \hTwo)$);

  \filldraw[fill=white] (M1t) -- (M2t) -- (M3t) -- (M4t) -- cycle;
  \filldraw[fill=white] (M4) -- (M1) -- (M1t) -- (M4t) -- cycle;
  \filldraw[fill=white] (M1) -- (M2) -- (M2t) -- (M1t) -- cycle;

  % ================= Layer 3 (Top cube) =================
  \coordinate (N1) at ($(M4t)!0.5!(M1t)$);
  \coordinate (N2) at ($(M1t)!0.5!(M2t)$);
  \coordinate (N3) at ($(M2t)!0.5!(M3t)$);
  \coordinate (N4) at ($(M3t)!0.5!(M4t)$);

  \def\hThree{0.8}
  \coordinate (N1t) at ($(N1) + (0, \hThree)$);
  \coordinate (N2t) at ($(N2) + (0, \hThree)$);
  \coordinate (N3t) at ($(N3) + (0, \hThree)$);
  \coordinate (N4t) at ($(N4) + (0, \hThree)$);

  \filldraw[fill=white] (N1t) -- (N2t) -- (N3t) -- (N4t) -- cycle;
  \filldraw[fill=white] (N1) -- (N2) -- (N2t) -- (N1t) -- cycle;
  \filldraw[fill=white] (N2) -- (N3) -- (N3t) -- (N2t) -- cycle;
\end{tikzpicture}
\end{center}
\nopagebreak[4]
\begin{gkopts}(4)
\task A．$4$
\task B．$5$
\task C．$6$
\task D．$7$
\end{gkopts}

\Needspace{4\baselineskip}
\textbf{解析}

设塔形几何体自下而上共有 $n$ 个正方体，记各层正方体棱长为 $a_k$（$k = 1, 2, \dots, n$）．
最底层正方体棱长 $a_1 = 2$，底面面积 $S_1 = a_1^2 = 4$．
由于上层正方体下底面的四个顶点是下层正方体上底面各边的中点，
\Needspace{5\baselineskip}
由平面几何性质，边中点连线所围成的正方形面积为原正方形面积的一半，即：
\nopagebreak[4]
\[
S_k = \frac{1}{2}S_{k-1} = \left(\frac{1}{2}\right)^{k-1} S_1 = 4 \times \left(\frac{1}{2}\right)^{k-1}.
\]
计算该塔形几何体的总表面积（含最底层底面面积）：
\Needspace{5\baselineskip}
1. 侧面积部分：每个正方体有 4 个侧面，各层正方体侧面积之和为：
\nopagebreak[4]
\[
S_{\text{\text{侧}}} = \sum_{k=1}^n 4 S_k = 4 \sum_{k=1}^n 4 \times \left(\frac{1}{2}\right)^{k-1} = 16 \sum_{k=1}^n \left(\frac{1}{2}\right)^{k-1}.
\]
2. 水平暴露表面部分：
最底层正方体的底面面积为 $S_1$；
自上向下俯视，各层正方体暴露的上表面投影刚好覆盖整个最底层正方形，其面积和为 $S_1$；
因此水平暴露表面总面积为 $S_1 + S_1 = 2S_1 = 8$．
\Needspace{5\baselineskip}
总表面积为：
\nopagebreak[4]
\[
S_{\text{\text{总}}} = 8 + 16 \times \frac{1 - (1/2)^n}{1 - 1/2} = 8 + 32\left(1 - \frac{1}{2^n}\right) = 40 - \frac{32}{2^n}.
\]
\Needspace{5\baselineskip}
依题意，$S_{\text{\text{总}}} > 39$：
\nopagebreak[4]
\[
40 - \frac{32}{2^n} > 39 \iff \frac{32}{2^n} < 1 \iff 2^n > 32 = 2^5.
\]
因为 $n$ 为正整数，所以 $n \geqslant 6$．
即塔形中正方体的个数至少是 6．
故选 C．

\gkanswer{答案：C}

\end{answer}
\begin{answer}
\gksolhead{11}{}

\textbf{题目}：

若集合 $A = \{x \in \mathbf{R} \mid x^2 - 4x + 3 < 0\}$，集合 $B = \{x \in \mathbf{R} \mid (x - 2)(x - 5) < 0\}$，则 $A \cap B =$ \underline{\hspace{2cm}}．

\Needspace{4\baselineskip}
\textbf{解析}

\Needspace{5\baselineskip}
解不等式 $x^2 - 4x + 3 < 0$：
\nopagebreak[4]
\[
(x - 1)(x - 3) < 0 \iff 1 < x < 3,
\]
故集合 $A = (1, 3)$．
集合 $B = \{x \in \mathbf{R} \mid (x - 2)(x - 5) < 0\} = (2, 5)$．
\Needspace{5\baselineskip}
取交集得：
\nopagebreak[4]
\[
A \cap B = (1, 3) \cap (2, 5) = (2, 3).
\]

\gkanswer{答案：$(2, 3)$（或 $\{x \mid 2 < x < 3\}$）}

\end{answer}
\begin{answer}
\gksolhead{12}{}

\textbf{题目}：

曲线 $y = x^3$ 在点 $(1, 1)$ 处的切线与 $x$ 轴、直线 $x = 2$ 所围成的三角形的面积为 \underline{\hspace{2cm}}．

\Needspace{4\baselineskip}
\textbf{解析}

求函数 $y = x^3$ 的导数：$y' = 3x^2$．
在点 $(1, 1)$ 处的切线斜率 $k = y'|_{x=1} = 3$．
切线方程为 $y - 1 = 3(x - 1)$，即 $y = 3x - 2$．
求切线与坐标轴及直线的交点：
令 $y = 0$，得 $3x - 2 = 0 \implies x = \dfrac{2}{3}$，交点为 $\left(\dfrac{2}{3}, 0\right)$；
直线 $x = 2$ 与 $x$ 轴交于 $(2, 0)$；
切线与直线 $x = 2$ 交于 $(2, 4)$．
\Needspace{5\baselineskip}
该三角形在 $x$ 轴上的底边长为：
\nopagebreak[4]
\[
2 - \frac{2}{3} = \frac{4}{3},
\]
底边上的高为交点的纵坐标 $h = 4$．
\Needspace{5\baselineskip}
所以三角形的面积为：
\nopagebreak[4]
\[
S = \frac{1}{2} \times \frac{4}{3} \times 4 = \frac{8}{3}.
\]

\gkanswer{答案：$\dfrac{8}{3}$}

\end{answer}
\begin{answer}
\gksolhead{13}{}

\textbf{题目}：

已知 $\alpha$、$\beta$ 均为锐角，且 $\cos(\alpha + \beta) = \sin(\alpha - \beta)$，则 $\tan\alpha =$ \underline{\hspace{2cm}}．

\Needspace{4\baselineskip}
\textbf{解析}

\Needspace{5\baselineskip}
由两角和与差的三角公式展开已知等式：
\nopagebreak[4]
\[
\cos(\alpha + \beta) = \cos\alpha\cos\beta - \sin\alpha\sin\beta,
\]
\nopagebreak[4]
\[
\sin(\alpha - \beta) = \sin\alpha\cos\beta - \cos\alpha\sin\beta.
\]
\Needspace{5\baselineskip}
因为 $\cos(\alpha + \beta) = \sin(\alpha - \beta)$，所以：
\nopagebreak[4]
\[
\cos\alpha\cos\beta - \sin\alpha\sin\beta = \sin\alpha\cos\beta - \cos\alpha\sin\beta.
\]
\Needspace{5\baselineskip}
移项并整理：
\nopagebreak[4]
\[
\cos\alpha(\cos\beta + \sin\beta) = \sin\alpha(\cos\beta + \sin\beta).
\]
因为 $\beta$ 为锐角，所以 $\cos\beta > 0, \sin\beta > 0 \implies \cos\beta + \sin\beta > 0$．
\Needspace{5\baselineskip}
两边同时除以 $\cos\beta + \sin\beta$，得：
\nopagebreak[4]
\[
\cos\alpha = \sin\alpha.
\]
因为 $\alpha$ 为锐角，两边同除以 $\cos\alpha$，得 $\tan\alpha = 1$．

\gkanswer{答案：$1$}

\end{answer}
\begin{answer}
\gksolhead{14}{}

\textbf{题目}：

若 $x^2 + y^2 = 4$，则 $x - y$ 的最大值是 \underline{\hspace{2cm}}．

\Needspace{4\baselineskip}
\textbf{解析}

方法一（三角代换法）：
由 $x^2 + y^2 = 4$，设 $x = 2\cos\theta, y = 2\sin\theta$（$\theta \in [0, 2\pi)$）．
\Needspace{5\baselineskip}
则：
\nopagebreak[4]
\[
x - y = 2\cos\theta - 2\sin\theta = 2\sqrt{2}\cos\left(\theta + \frac{\pi}{4}\right).
\]
因为 $\cos\left(\theta + \dfrac{\pi}{4}\right) \leqslant 1$，当且仅当 $\theta = -\dfrac{\pi}{4}$ 时取等号，
所以 $x - y$ 的最大值是 $2\sqrt{2}$．

方法二（柯西不等式）：
\Needspace{5\baselineskip}
由柯西不等式：
\nopagebreak[4]
\[
(x - y)^2 = (1 \cdot x + (-1) \cdot y)^2 \leqslant (1^2 + (-1)^2)(x^2 + y^2) = 2 \times 4 = 8.
\]
故 $-\sqrt{8} \leqslant x - y \leqslant \sqrt{8} = 2\sqrt{2}$，最大值为 $2\sqrt{2}$．

\gkanswer{答案：$2\sqrt{2}$}

\end{answer}
\begin{answer}
\gksolhead{15}{}

\textbf{题目}：

若 10 把钥匙中只有 2 把能打开某锁，则从中任取 2 把能将该锁打开的概率为 \underline{\hspace{2cm}}．

\Needspace{4\baselineskip}
\textbf{解析}

\Needspace{5\baselineskip}
从 10 把钥匙中任取 2 把，基本事件总数为：
\nopagebreak[4]
\[
\mathrm{C}_{10}^2 = \frac{10 \times 9}{2} = 45.
\]
设事件 $A$ 为“任取 2 把能将该锁打开”，其对立事件 $\bar{A}$ 为“任取的 2 把均不能打开该锁”．
因为 10 把钥匙中只有 2 把能开锁，不能开锁的钥匙有 8 把，
\Needspace{5\baselineskip}
从这 8 把不能开锁的钥匙中任取 2 把的取法共有：
\nopagebreak[4]
\[
\mathrm{C}_8^2 = \frac{8 \times 7}{2} = 28.
\]
所以对立事件的概率为 $P(\bar{A}) = \dfrac{28}{45}$．
\Needspace{5\baselineskip}
由对立事件概率公式得：
\nopagebreak[4]
\[
P(A) = 1 - P(\bar{A}) = 1 - \frac{28}{45} = \frac{17}{45}.
\]

\gkanswer{答案：$\dfrac{17}{45}$}

\end{answer}
\begin{answer}
\gksolhead{16}{}

\textbf{题目}：

已知 $A\left(-\dfrac{1}{2}, 0\right)$，$B$ 是圆 $F: \left(x - \dfrac{1}{2}\right)^2 + y^2 = 4$（$F$ 为圆心）上一动点，线段 $AB$ 的垂直平分线交 $BF$ 于 $P$，则动点 $P$ 的轨迹方程为 \underline{\hspace{2cm}}．

\Needspace{4\baselineskip}
\textbf{解析}

圆 $F: \left(x - \dfrac{1}{2}\right)^2 + y^2 = 4$ 的圆心为 $F\left(\dfrac{1}{2}, 0\right)$，半径为 $R = 2$．
动点 $B$ 在圆 $F$ 上，所以 $|BF| = 2$．
已知点 $A\left(-\dfrac{1}{2}, 0\right)$，点 $P$ 在线段 $AB$ 的垂直平分线上，由垂直平分线性质知 $|PA| = |PB|$．
因为点 $P$ 是线段 $AB$ 的垂直平分线与半径 $BF$ 的交点，且 $|AF| = 1 < 2$（点 $A$ 在圆内），
\Needspace{5\baselineskip}
所以点 $P$ 在线段 $BF$ 上，满足：
\nopagebreak[4]
\[
|PA| + |PF| = |PB| + |PF| = |BF| = 2.
\]
由椭圆的第一定义，动点 $P$ 到两定点 $A\left(-\dfrac{1}{2}, 0\right)$ 和 $F\left(\dfrac{1}{2}, 0\right)$ 的距离之和为定值 $2$，且 $2 > |AF| = 1$，
因此动点 $P$ 的轨迹是以 $A, F$ 为焦点的椭圆．
该椭圆的中心在原点 $(0, 0)$，焦点在 $x$ 轴上：
半长轴 $a = 1$，半焦距 $c = \dfrac{1}{2}$，
半短轴 $b^2 = a^2 - c^2 = 1 - \dfrac{1}{4} = \dfrac{3}{4}$．
\Needspace{5\baselineskip}
故动点 $P$ 的轨迹方程为：
\nopagebreak[4]
\[
x^2 + \frac{4y^2}{3} = 1.
\]

\gkanswer{答案：$x^2 + \dfrac{4y^2}{3} = 1$（或 $\dfrac{x^2}{1} + \dfrac{y^2}{\frac{3}{4}} = 1$）}

\end{answer}
\begin{answer}
\gksolhead{17}{}

\textbf{题目}：

若函数 $f(x) = \dfrac{1 + \cos 2x}{2\sin\left(\dfrac{\pi}{2} - x\right)} + \sin x + a^2 \sin\left(x + \dfrac{\pi}{4}\right)$ 的最大值为 $\sqrt{2} + 3$，试确定常数 $a$ 的值．

\Needspace{4\baselineskip}
\textbf{解析}

由诱导公式和二倍角公式化简函数解析式：
分母 $2\sin\left(\dfrac{\pi}{2} - x\right) = 2\cos x$（定义域要求 $\cos x \neq 0$）．
分子 $1 + \cos 2x = 2\cos^2 x$．
\Needspace{5\baselineskip}
因此前一项为：
\nopagebreak[4]
\[
\frac{1 + \cos 2x}{2\sin\left(\dfrac{\pi}{2} - x\right)} = \frac{2\cos^2 x}{2\cos x} = \cos x.
\]
\Needspace{5\baselineskip}
代入 $f(x)$ 得：
\nopagebreak[4]
\[
f(x) = \cos x + \sin x + a^2 \sin\left(x + \frac{\pi}{4}\right).
\]
\Needspace{5\baselineskip}
将 $\cos x + \sin x$ 辅助角化简：
\nopagebreak[4]
\[
\cos x + \sin x = \sqrt{2}\left(\frac{\sqrt{2}}{2}\cos x + \frac{\sqrt{2}}{2}\sin x\right) = \sqrt{2}\sin\left(x + \frac{\pi}{4}\right).
\]
\Needspace{5\baselineskip}
所以：
\nopagebreak[4]
\[
f(x) = (\sqrt{2} + a^2)\sin\left(x + \frac{\pi}{4}\right).
\]
因为 $a \in \mathbf{R}$，所以 $a^2 \geqslant 0 \implies \sqrt{2} + a^2 > 0$．
当 $\sin\left(x + \dfrac{\pi}{4}\right) = 1$ 时（此时 $x + \dfrac{\pi}{4} = \dfrac{\pi}{2} + 2k\pi \implies x = \dfrac{\pi}{4} + 2k\pi$，满足 $\cos x \neq 0$），
\Needspace{5\baselineskip}
函数取得最大值：
\nopagebreak[4]
\[
f_{\max} = \sqrt{2} + a^2.
\]
\Needspace{5\baselineskip}
已知最大值为 $\sqrt{2} + 3$，故：
\nopagebreak[4]
\[
\sqrt{2} + a^2 = \sqrt{2} + 3 \implies a^2 = 3 \implies a = \pm\sqrt{3}.
\]

\gkanswer{答案：$a = \pm\sqrt{3}$。}

\end{answer}
\begin{answer}
\gksolhead{18}{}

\textbf{题目}：

加工某种零件需经过三道工序，设第一、二、三道工序的合格率分别为 $\dfrac{9}{10}$、$\dfrac{8}{9}$、$\dfrac{7}{8}$，且各道工序互不影响．

(1) 求该种零件的合格率；

(2) 从该种零件中任取 3 件，求恰好取到一件合格品的概率和至少取到一件合格品的概率．

\Needspace{4\baselineskip}
\textbf{解析}

\gksubpart{(1)}

设事件 $A_1, A_2, A_3$ 分别表示第一、二、三道工序合格．
由题意，各道工序互不影响，即事件 $A_1, A_2, A_3$ 相互独立，
且 $P(A_1) = \dfrac{9}{10}$，$P(A_2) = \dfrac{8}{9}$，$P(A_3) = \dfrac{7}{8}$．
\Needspace{5\baselineskip}
零件合格当且仅当三道工序都合格，记零件合格的事件为 $A$：
\nopagebreak[4]
\[
P(A) = P(A_1 A_2 A_3) = P(A_1) P(A_2) P(A_3) = \frac{9}{10} \times \frac{8}{9} \times \frac{7}{8} = \frac{7}{10}.
\]
答：该种零件的合格率为 $\dfrac{7}{10}$（或 $0.7$）．

\gksubpart{(2)}

从该种零件中任取 3 件，每件合格的概率为 $p = \dfrac{7}{10}$，不合格的概率为 $1 - p = \dfrac{3}{10}$．
各件取样相互独立，服从二项分布．
\Needspace{5\baselineskip}
1. 恰好取到一件合格品的概率为：
\nopagebreak[4]
\[
P_1 = \mathrm{C}_3^1 p^1 (1 - p)^2 = 3 \times \frac{7}{10} \times \left(\frac{3}{10}\right)^2 = \frac{189}{1000} = 0.189.
\]
\Needspace{5\baselineskip}
2. 取出的 3 件全为不合格品的概率为：
\nopagebreak[4]
\[
P_0 = (1 - p)^3 = \left(\frac{3}{10}\right)^3 = \frac{27}{1000}.
\]
\Needspace{5\baselineskip}
故至少取到一件合格品的概率为：
\nopagebreak[4]
\[
P_2 = 1 - P_0 = 1 - \frac{27}{1000} = \frac{973}{1000} = 0.973.
\]

\gkanswer{答案（1）：$\dfrac{7}{10}$（或 $0.7$）。}

\gkanswer{答案（2）：恰好取到一件合格品的概率为 $\dfrac{189}{1000}$；至少取到一件合格品的概率为 $\dfrac{973}{1000}$。}

\end{answer}
\begin{answer}
\gksolhead{19}{}

\textbf{题目}：

设函数 $f(x) = 2x^3 - 3(a + 1)x^2 + 6ax + 8$，其中 $a \in \mathbf{R}$．

(1) 若 $f(x)$ 在 $x = 3$ 处取得极值，求常数 $a$ 的值；

(2) 若 $f(x)$ 在 $(-\infty, 0)$ 上为增函数，求 $a$ 的取值范围．

\Needspace{4\baselineskip}
\textbf{解析}

\Needspace{5\baselineskip}
对函数 $f(x) = 2x^3 - 3(a + 1)x^2 + 6ax + 8$ 求导：
\nopagebreak[4]
\[
f'(x) = 6x^2 - 6(a + 1)x + 6a = 6\left[x^2 - (a + 1)x + a\right] = 6(x - 1)(x - a).
\]

\gksubpart{(1)}

若 $f(x)$ 在 $x = 3$ 处取得极值，则必有 $f'(3) = 0$．
\Needspace{5\baselineskip}
代入得：
\nopagebreak[4]
\[
f'(3) = 6(3 - 1)(3 - a) = 12(3 - a) = 0 \implies a = 3.
\]
检验：当 $a = 3$ 时，$f'(x) = 6(x - 1)(x - 3)$．
当 $x \in (1, 3)$ 时，$f'(x) < 0$；当 $x \in (3, +\infty)$ 时，$f'(x) > 0$．
因此 $x = 3$ 是 $f(x)$ 的极小值点，符合题意．
故 $a = 3$．

\gksubpart{(2)}

若 $f(x)$ 在 $(-\infty, 0)$ 上为增函数，则对于任意 $x \in (-\infty, 0)$，恒有 $f'(x) \geqslant 0$．
\Needspace{5\baselineskip}
即：
\nopagebreak[4]
\[
(x - 1)(x - a) \geqslant 0 \quad (\forall x \in (-\infty, 0)).
\]
\Needspace{5\baselineskip}
当 $x < 0$ 时，$x - 1 < 0$ 恒成立，两边同除以 $x - 1$ 得：
\nopagebreak[4]
\[
x - a \leqslant 0 \iff a \geqslant x \quad (\forall x \in (-\infty, 0)).
\]
要使 $a \geqslant x$ 对任意 $x < 0$ 恒成立，只需 $a \geqslant 0$．
检验 $a = 0$：此时 $f'(x) = 6x(x - 1)$，当 $x \in (-\infty, 0)$ 时，$x < 0, x - 1 < 0 \implies f'(x) > 0$，为增函数．
所以实数 $a$ 的取值范围是 $[0, +\infty)$．

\gkanswer{答案（1）：$a = 3$。}

\gkanswer{答案（2）：$[0, +\infty)$（或 $a \geqslant 0$）。}

\end{answer}
\begin{answer}
\gksolhead{20}{}

\textbf{题目}：

如图，在四棱锥 $P - ABCD$ 中，底面 $ABCD$ 为矩形，$PD \perp$ 底面 $ABCD$，$E$ 是 $AB$ 上一点，$PE \perp EC$．已知 $PD = \sqrt{2}, CD = 2, AE = \dfrac{1}{2}$，求：
\nopagebreak[4]
\begin{center}\begin{tikzpicture}[>=stealth, scale=1.3, line join=round, line cap=round, thick]
  % Coordinates
  \coordinate (A) at (0, 0);
  \coordinate (D) at (0.8, 0.7);
  \coordinate (P) at (0.8, 2.8);
  \coordinate (E) at (1.4, 0);
  \coordinate (B) at (3.6, 0);
  \coordinate (C) at ($(B) + (D) - (A)$);

  % Dashed lines (hidden / interior)
  \draw[dashed] (P) -- (D);
  \draw[dashed] (D) -- (A);
  \draw[dashed] (D) -- (C);
  \draw[dashed] (D) -- (E);
  \draw[dashed] (E) -- (C);

  % Solid lines (outer / visible)
  \draw (P) -- (A);
  \draw (P) -- (E);
  \draw (P) -- (B);
  \draw (P) -- (C);
  \draw (A) -- (B) -- (C);

  % Labels
  \node[above] at (P) {$P$};
  \node[above right] at (D) {$D$};
  \node[below left] at (A) {$A$};
  \node[below] at (E) {$E$};
  \node[below right] at (B) {$B$};
  \node[above right] at (C) {$C$};
\end{tikzpicture}
\end{center}
\nopagebreak[4]
(1) 异面直线 $PD$ 与 $EC$ 的距离；

(2) 二面角 $E - PC - D$ 的大小．

\Needspace{4\baselineskip}
\textbf{解析}

\gksubpart{(1)}

因为底面 $ABCD$ 为矩形，所以 $AD \perp AB, CD \perp AD, AB \parallel CD, AB = CD = 2$．
已知 $E$ 是 $AB$ 上一点，$AE = \dfrac{1}{2}$，所以 $EB = AB - AE = 2 - \dfrac{1}{2} = \dfrac{3}{2}$．
因为 $PD \perp$ 底面 $ABCD$，且 $DE \subset$ 底面 $ABCD$，所以 $PD \perp DE$．
$DE$ 是斜线 $PE$ 在底面 $ABCD$ 上的射影．
又已知 $PE \perp EC$，由三垂线定理的逆定理，得 $DE \perp EC$．
在直角三角形 $DAE$ 中，$DE^2 = AD^2 + AE^2 = AD^2 + \dfrac{1}{4}$．
在直角三角形 $EBC$ 中，$EC^2 = BC^2 + EB^2 = AD^2 + \dfrac{9}{4}$．
在直角三角形 $DEC$ 中，$DE^2 + EC^2 = CD^2 = 2^2 = 4$．
\Needspace{5\baselineskip}
代入得：
\nopagebreak[4]
\[
\left(AD^2 + \frac{1}{4}\right) + \left(AD^2 + \frac{9}{4}\right) = 4 \implies 2AD^2 + \frac{5}{2} = 4 \implies 2AD^2 = \frac{3}{2} \implies AD^2 = \frac{3}{4}.
\]
解得 $AD = \dfrac{\sqrt{3}}{2}$．
从而 $DE = \sqrt{AD^2 + AE^2} = \sqrt{\dfrac{3}{4} + \dfrac{1}{4}} = 1$．
因为 $PD \perp$ 底面 $ABCD$，所以 $PD \perp DE$；
又 $DE \perp EC$，且 $PD$ 与 $EC$ 互为异面直线，
所以线段 $DE$ 是异面直线 $PD$ 与 $EC$ 的公垂线段．
因此，异面直线 $PD$ 与 $EC$ 的距离为 $DE = 1$．

\gksubpart{(2)}

以 $D$ 为坐标原点，$DA$ 所在直线为 $x$ 轴，$DC$ 所在直线为 $y$ 轴，$DP$ 所在直线为 $z$ 轴，建立空间直角坐标系 $D - xyz$．
根据已知及第 (1) 问的结论，各点坐标为：
$D(0, 0, 0)$，$A\left(\dfrac{\sqrt{3}}{2}, 0, 0\right)$，$C(0, 2, 0)$，$P(0, 0, \sqrt{2})$．
因为点 $E$ 在 $AB$ 上且 $AE = \dfrac{1}{2}$，所以 $E\left(\dfrac{\sqrt{3}}{2}, \dfrac{1}{2}, 0\right)$．
因为 $D, C, P$ 均在 $yOz$ 平面上，所以平面 $PCD$ 即为 $yOz$ 平面，
其法向量可取为 $\vec{n}_1 = (1, 0, 0)$．
\Needspace{5\baselineskip}
求平面 $PCE$ 的法向量 $\vec{n}_2 = (x, y, z)$：
\nopagebreak[4]
\[
\vec{PC} = (0, 2, -\sqrt{2}), \quad \vec{EC} = \left(-\frac{\sqrt{3}}{2}, \frac{3}{2}, 0\right).
\]
\Needspace{5\baselineskip}
由 $\vec{n}_2 \perp \vec{PC}$ 及 $\vec{n}_2 \perp \vec{EC}$ 得：
\nopagebreak[4]
\[
\begin{cases}
2y - \sqrt{2}z = 0, \\
-\dfrac{\sqrt{3}}{2}x + \dfrac{3}{2}y = 0.
\end{cases}
\]
解得 $z = \sqrt{2}y$，$x = \sqrt{3}y$．
取 $y = 1$，得 $\vec{n}_2 = (\sqrt{3}, 1, \sqrt{2})$．
\Needspace{5\baselineskip}
其模长为：
\nopagebreak[4]
\[
|\vec{n}_2| = \sqrt{(\sqrt{3})^2 + 1^2 + (\sqrt{2})^2} = \sqrt{6}.
\]
\Needspace{5\baselineskip}
设二面角 $E - PC - D$ 的大小为 $\theta$（$\theta$ 为锐角），则：
\nopagebreak[4]
\[
\cos\theta = \frac{|\vec{n}_1 \cdot \vec{n}_2|}{|\vec{n}_1||\vec{n}_2|} = \frac{\sqrt{3}}{1 \times \sqrt{6}} = \frac{1}{\sqrt{2}} = \frac{\sqrt{2}}{2}.
\]
所以 $\theta = \dfrac{\pi}{4}$．
答：二面角 $E - PC - D$ 的大小为 $\dfrac{\pi}{4}$（或 $45^\circ$）．

\gkanswer{答案（1）：$1$。}

\gkanswer{答案（2）：$\dfrac{\pi}{4}$（或 $45^\circ$）。}

\end{answer}
\begin{answer}
\gksolhead{21}{}

\textbf{题目}：

已知中心在原点的双曲线 $C$ 的右焦点为 $(2, 0)$，右顶点为 $(\sqrt{3}, 0)$．

(1) 求双曲线 $C$ 的方程；

(2) 若直线 $l: y = kx + \sqrt{2}$ 与双曲线 $C$ 恒有两个不同的交点 $A$ 和 $B$，且 $\overrightarrow{OA} \cdot \overrightarrow{OB} > 2$（其中 $O$ 为原点）．求 $k$ 的取值范围．

\Needspace{4\baselineskip}
\textbf{解析}

\gksubpart{(1)}

由双曲线中心在原点，焦点在 $x$ 轴上，
右焦点为 $(2, 0) \implies c = 2$，
右顶点为 $(\sqrt{3}, 0) \implies a = \sqrt{3}$，
所以 $a^2 = 3$．
\Needspace{5\baselineskip}
由双曲线几何性质 $b^2 = c^2 - a^2$，得：
\nopagebreak[4]
\[
b^2 = 2^2 - (\sqrt{3})^2 = 4 - 3 = 1.
\]
所以双曲线 $C$ 的方程为 $\dfrac{x^2}{3} - y^2 = 1$．

\gksubpart{(2)}

\Needspace{5\baselineskip}
将直线方程 $y = kx + \sqrt{2}$ 代入双曲线方程 $\dfrac{x^2}{3} - y^2 = 1$，消去 $y$：
\nopagebreak[4]
\[
x^2 - 3(kx + \sqrt{2})^2 = 3 \implies (1 - 3k^2)x^2 - 6\sqrt{2}kx - 9 = 0.
\]
因为直线与双曲线有两个不同的交点 $A(x_1, y_1), B(x_2, y_2)$，所以：
\Needspace{5\baselineskip}
1. 二次项系数不为 0：
\nopagebreak[4]
\[
1 - 3k^2 \neq 0 \iff k^2 \neq \frac{1}{3}.
\]
\Needspace{5\baselineskip}
2. 判别式 $\Delta > 0$：
\nopagebreak[4]
\[
\Delta = (-6\sqrt{2}k)^2 - 4(1 - 3k^2)(-9) = 72k^2 + 36(1 - 3k^2) = 36(1 - k^2) > 0 \iff k^2 < 1.
\]
\Needspace{5\baselineskip}
由韦达定理：
\nopagebreak[4]
\[
x_1 + x_2 = \frac{6\sqrt{2}k}{1 - 3k^2}, \quad x_1 x_2 = \frac{-9}{1 - 3k^2}.
\]
\Needspace{5\baselineskip}
计算 $\overrightarrow{OA} \cdot \overrightarrow{OB}$：
\nopagebreak[4]
\[
\overrightarrow{OA} \cdot \overrightarrow{OB} = x_1 x_2 + y_1 y_2.
\]
\Needspace{5\baselineskip}
因为 $y_1 = kx_1 + \sqrt{2}, y_2 = kx_2 + \sqrt{2}$，所以：
\nopagebreak[4]
\[
y_1 y_2 = k^2 x_1 x_2 + \sqrt{2}k(x_1 + x_2) + 2.
\]
\Needspace{5\baselineskip}
故：
\nopagebreak[4]
\[
\overrightarrow{OA} \cdot \overrightarrow{OB} = (1 + k^2)x_1 x_2 + \sqrt{2}k(x_1 + x_2) + 2.
\]
\Needspace{5\baselineskip}
代入韦达定理：
\nopagebreak[4]
\[
\overrightarrow{OA} \cdot \overrightarrow{OB} = (1 + k^2)\left(\frac{-9}{1 - 3k^2}\right) + \sqrt{2}k\left(\frac{6\sqrt{2}k}{1 - 3k^2}\right) + 2.
\]
\Needspace{5\baselineskip}
化简分子：
\nopagebreak[4]
\[
\overrightarrow{OA} \cdot \overrightarrow{OB} = \frac{-9 - 9k^2 + 12k^2 + 2(1 - 3k^2)}{1 - 3k^2} = \frac{-3k^2 - 7}{1 - 3k^2}.
\]
\Needspace{5\baselineskip}
依题意 $\overrightarrow{OA} \cdot \overrightarrow{OB} > 2$：
\nopagebreak[4]
\[
\frac{-3k^2 - 7}{1 - 3k^2} > 2 \iff \frac{-3k^2 - 7 - 2(1 - 3k^2)}{1 - 3k^2} > 0 \iff \frac{3k^2 - 9}{1 - 3k^2} > 0.
\]
\Needspace{5\baselineskip}
即：
\nopagebreak[4]
\[
\frac{3 - k^2}{3k^2 - 1} > 0.
\]
因为由前面判别式条件 $k^2 < 1$，分子 $3 - k^2 > 3 - 1 = 2 > 0$ 恒成立，
\Needspace{5\baselineskip}
所以只需分母大于 0：
\nopagebreak[4]
\[
3k^2 - 1 > 0 \iff k^2 > \frac{1}{3}.
\]
\Needspace{5\baselineskip}
结合 $k^2 < 1$，得：
\nopagebreak[4]
\[
\frac{1}{3} < k^2 < 1.
\]
解得 $k \in \left(-1, -\dfrac{\sqrt{3}}{3}\right) \cup \left(\dfrac{\sqrt{3}}{3}, 1\right)$．

\gkanswer{答案（1）：$\dfrac{x^2}{3} - y^2 = 1$。}

\gkanswer{答案（2）：$\left(-1, -\dfrac{\sqrt{3}}{3}\right) \cup \left(\dfrac{\sqrt{3}}{3}, 1\right)$。}

\end{answer}
\begin{answer}
\gksolhead{22}{}

\textbf{题目}：

数列 $\{a_n\}$ 满足 $a_1 = 1$ 且 $8a_{n+1} a_n - 16a_{n+1} + 2a_n + 5 = 0$（$n \geqslant 1$）．记 $b_n = \dfrac{1}{a_n - \dfrac{1}{2}}$（$n \geqslant 1$）．

(1) 求 $b_1$、$b_2$、$b_3$、$b_4$ 的值；

(2) 求数列 $\{b_n\}$ 的通项公式及数列 $\{a_n b_n\}$ 的前 $n$ 项和 $S_n$．

\Needspace{4\baselineskip}
\textbf{解析}

\gksubpart{(1)}

\Needspace{5\baselineskip}
已知 $a_1 = 1$，由 $b_n = \dfrac{1}{a_n - 1/2}$ 得：
\nopagebreak[4]
\[
b_1 = \frac{1}{a_1 - 1/2} = \frac{1}{1 - 1/2} = 2.
\]
\Needspace{5\baselineskip}
由已知递推关系 $8a_{n+1} a_n - 16a_{n+1} + 2a_n + 5 = 0$（$n \geqslant 1$），得：
\nopagebreak[4]
\[
8a_{n+1}(a_n - 2) = -(2a_n + 5) \implies a_{n+1} = \frac{2a_n + 5}{16 - 8a_n}.
\]
\Needspace{5\baselineskip}
从而：
\nopagebreak[4]
\[
a_{n+1} - \frac{1}{2} = \frac{2a_n + 5}{16 - 8a_n} - \frac{1}{2} = \frac{2a_n + 5 - (8 - 4a_n)}{16 - 8a_n} = \frac{6a_n - 3}{16 - 8a_n} = \frac{6(a_n - 1/2)}{16 - 8a_n}.
\]
\Needspace{5\baselineskip}
两边取倒数，得：
\nopagebreak[4]
\[
b_{n+1} = \frac{1}{a_{n+1} - 1/2} = \frac{16 - 8a_n}{6(a_n - 1/2)} = \frac{-8(a_n - 1/2) + 12}{6(a_n - 1/2)} = 2b_n - \frac{4}{3}.
\]
\Needspace{5\baselineskip}
由 $b_1 = 2$，依次递推计算：
\nopagebreak[4]
\[
b_2 = 2b_1 - \frac{4}{3} = 2 \times 2 - \frac{4}{3} = \frac{8}{3},
\]
\nopagebreak[4]
\[
b_3 = 2b_2 - \frac{4}{3} = 2 \times \frac{8}{3} - \frac{4}{3} = 4,
\]
\nopagebreak[4]
\[
b_4 = 2b_3 - \frac{4}{3} = 2 \times 4 - \frac{4}{3} = \frac{20}{3}.
\]

\gksubpart{(2)}

\Needspace{5\baselineskip}
由 $b_{n+1} = 2b_n - \dfrac{4}{3}$，两边同减去 $\dfrac{4}{3}$：
\nopagebreak[4]
\[
b_{n+1} - \frac{4}{3} = 2\left(b_n - \frac{4}{3}\right).
\]
因为 $b_1 - \dfrac{4}{3} = 2 - \dfrac{4}{3} = \dfrac{2}{3} \neq 0$，
所以数列 $\left\{b_n - \dfrac{4}{3}\right\}$ 是以 $\dfrac{2}{3}$ 为首项、2 为公比的等比数列．
\Needspace{5\baselineskip}
通项公式为：
\nopagebreak[4]
\[
b_n - \frac{4}{3} = \frac{2}{3} \times 2^{n-1} = \frac{2^n}{3} \implies b_n = \frac{2^n + 4}{3}.
\]
因为 $b_n = \dfrac{1}{a_n - 1/2}$，所以 $a_n = \dfrac{1}{b_n} + \dfrac{1}{2}$．
\Needspace{5\baselineskip}
从而：
\nopagebreak[4]
\[
a_n b_n = \left(\frac{1}{b_n} + \frac{1}{2}\right) b_n = 1 + \frac{1}{2}b_n = 1 + \frac{1}{2} \times \frac{2^n + 4}{3} = \frac{2^{n-1} + 5}{3}.
\]
\Needspace{5\baselineskip}
求数列 $\{a_n b_n\}$ 的前 $n$ 项和 $S_n$：
\nopagebreak[4]
\[
S_n = \sum_{k=1}^n a_k b_k = \frac{1}{3} \sum_{k=1}^n 2^{k-1} + \frac{5n}{3}.
\]
\Needspace{5\baselineskip}
由等比数列求和公式：
\nopagebreak[4]
\[
S_n = \frac{1}{3} \times \frac{1 - 2^n}{1 - 2} + \frac{5n}{3} = \frac{2^n - 1}{3} + \frac{5n}{3} = \frac{2^n + 5n - 1}{3}.
\]

\gkanswer{答案（1）：$b_1 = 2$，$b_2 = \dfrac{8}{3}$，$b_3 = 4$，$b_4 = \dfrac{20}{3}$。}

\gkanswer{答案（2）：$b_n = \dfrac{2^n + 4}{3}$；$S_n = \dfrac{2^n + 5n - 1}{3}$。}

\end{answer}